% !TeX program = xelatex
% chktex-file 1
% chktex-file 44
% XeLaTeX-only version for Source Han Sans / Noto Sans CJK typography.
\documentclass[11pt,a4paper]{article}
\usepackage[margin=0.85in]{geometry}

\usepackage{iftex}
\ifXeTeX
\else
\PackageError{main}{This document requires XeLaTeX}{Please compile with xelatex.}
\fi
\usepackage{fontspec}
\usepackage{xeCJK}

\setmainfont{Helvetica Neue}
\setsansfont{Helvetica Neue}
\setmonofont{Menlo}
\setCJKmainfont{Hiragino Sans W3}[BoldFont=Hiragino Sans W6]
\setCJKsansfont{Hiragino Sans W3}[BoldFont=Hiragino Sans W6]
\setCJKmonofont{Hiragino Sans W3}[BoldFont=Hiragino Sans W6]
\renewcommand{\CJKfamilydefault}{\CJKsfdefault}

\usepackage{amsmath,amssymb}
\usepackage[table]{xcolor}
\usepackage{tikz-cd}
\usepackage{multicol}
\usepackage{tabularx,xltabular,needspace}
\usepackage{array}
\usepackage{enumitem}
\usepackage{titlesec}
\usepackage{tcolorbox}
\usepackage[bookmarks=false,hidelinks]{hyperref}
\renewcommand{\familydefault}{\sfdefault}

\setlength{\parindent}{0pt}
\setlength{\parskip}{0.45\baselineskip}
\renewcommand{\arraystretch}{1.18}
\setlength{\tabcolsep}{4.8pt}
\emergencystretch=2em
\clubpenalty=10000\widowpenalty=10000
\sloppy

\definecolor{MainBlue}{HTML}{3A3A3A}
\definecolor{DeepNavy}{HTML}{111111}
\definecolor{AccentTeal}{HTML}{555555}
\definecolor{AccentCoral}{HTML}{333333}
\definecolor{LightBlue}{HTML}{F5F5F5}
\definecolor{SoftGray}{HTML}{FAFAFA}
\definecolor{DarkGray}{HTML}{333333}
\definecolor{WarnYellow}{HTML}{EDEDED}

\newcommand{\overviewitem}[4]{%
  \noindent
  \begin{minipage}[t]{0.12\textwidth}
    \textbf{#1}\\[-0.5mm]
    {\scriptsize #4}
  \end{minipage}
  \hfill
  \begin{minipage}[t]{0.84\textwidth}
    {\scriptsize\textbf{Lecture:} #2}\\[-0.2mm]
    {\scriptsize\textbf{Practice:} #3}
  \end{minipage}
  \par\vspace{1.2mm}
  \hrule
  \vspace{1.2mm}
}

\tcbuselibrary{breakable}
\newenvironment{projectfiles}{
  \vspace{2.5mm}
  \begin{tcolorbox}[colback=SoftGray,colframe=DarkGray,arc=1.5mm,boxrule=0.45pt,left=2.5mm,right=2.5mm,top=2mm,bottom=1.5mm,breakable]
    {\normalsize\bfseries Workshop files}\\[-0.2mm]
    \begin{enumerate}[leftmargin=1.7em,itemsep=2pt,topsep=2pt,parsep=0pt]
      }{
    \end{enumerate}
  \end{tcolorbox}
}

\titleformat{\section}{\Large\bfseries\color{MainBlue}}{}{0pt}{}
\titleformat{\subsection}{\large\bfseries\color{DarkGray}}{}{0pt}{}
\renewcommand{\tabularxcolumn}[1]{m{#1}}
\newcolumntype{Y}{>{\raggedright\arraybackslash}X}
\newcolumntype{C}[1]{>{\centering\arraybackslash}m{#1}}
\newcolumntype{L}[1]{>{\raggedright\arraybackslash}m{#1}}
\newcommand{\checkbox}{\(\square\)}
\newcommand{\donebox}{\(\blacksquare\)}
\newcommand{\typebadge}[1]{\colorbox{LightBlue}{\textcolor{MainBlue}{\scriptsize\bfseries\strut #1}}}
\newcommand{\checks}[1]{\begin{itemize}[leftmargin=1.45em,itemsep=0pt,topsep=0pt,parsep=0pt,label=\(\square\)]#1\end{itemize}}
\newcommand{\daybox}[2]{%
  \begin{tcolorbox}[colback=SoftGray,colframe=DarkGray,arc=1mm,boxrule=0.45pt,left=2.5mm,right=2.5mm,top=2mm,bottom=2mm]
    \textbf{Objectives:} #1

    \vspace{1mm}
    \textbf{Outputs:} #2
  \end{tcolorbox}
}
\newcommand{\projectfile}[2]{\item {\ttfamily\small #1}\\[-0.3mm]{\footnotesize #2}}

\begin{document}
\newcommand{\tasktableheader}{%
  \rowcolor{MainBlue}
  \textcolor{white}{\textbf{Done}} &
  \textcolor{white}{\textbf{ID}} &
  \textcolor{white}{\textbf{Type}} &
  \textcolor{white}{\textbf{Task}} &
  \textcolor{white}{\textbf{Assessment criteria}} &
  \textcolor{white}{\textbf{Pts}} \\ \hline
}

\begin{center}
  {\LARGE\bfseries Python Basketball Analytics}\\[4pt]
  {\Large Five-Day Hackathon Schedule and Assessment Checklist}\\[10pt]
  \rule{0.9\textwidth}{0.6pt}\\[8pt]
  {\large \textcolor{DarkGray}{Colab | Roboflow | Homography | Detection | ByteTrack | Pose | Shooting Analysis}}
\end{center}
\vspace{2mm}

\begin{tcolorbox}[colback=SoftGray,colframe=DarkGray,arc=1mm,boxrule=0.45pt,left=3mm,right=3mm,top=2mm,bottom=2mm]
  {\large\bfseries Course overview}
  \vspace{2mm}\hrule\vspace{1.5mm}

\overviewitem{Day 1}{Colab, image coordinates, keypoint correspondences, homography and BEV projection.}{Match keypoints in a fixed image and a BEV court diagram, project a point, and prepare Roboflow keypoint and box annotations.}{20 pts}
  \overviewitem{Day 2}{Object detection, bounding boxes, confidence, bottom-center footpoints and BBOX-to-BEV.}{Start with manually specified boxes, examine sample detections or a trained YOLO26 model, then project multiple player boxes onto BEV.}{25 pts}
  \overviewitem{Day 3}{Multi-object tracking, IoU association, ByteTrack, track IDs, BEV paths and team clustering.}{Link detections into tracks and project them onto BEV. Use torso crops and HSV color features to form Team A and Team B.}{20 pts}
  \overviewitem{Day 4}{Close-up shooting footage, the provided shot, pose and ball models, shot events and release angles.}{Upload a shooting video through Google Drive or Colab. Estimate the shooter, body keypoints and ball path, then derive the release frame, joint angles and 2D launch angle.}{15 pts}
  \overviewitem{Day 5}{Research questions, system architecture, evaluation design, limitations and technical communication.}{Present a project proposal or demonstrate a completed basketball vision analysis project.}{20 pts}

  \noindent
  \begin{minipage}[t]{0.12\textwidth}\textbf{Bonus}\\[-0.5mm]{\scriptsize +10}\end{minipage}
  \hfill
  \begin{minipage}[t]{0.84\textwidth}{\scriptsize Improve a demo, analyze failures, refine a visualization, analyze another video, or submit specific course feedback.}\end{minipage}

  \vspace{2mm}\hrule\vspace{1.5mm}
  \hfill {\bfseries Total: 100 points (+10 bonus)}
\end{tcolorbox}

\clearpage

\section{Day 1: Colab, Keypoint Correspondences, Homography and Annotation}
\daybox{Become familiar with Colab, the course folders and image coordinates. Establish correspondences between a fixed camera image and a BEV diagram, then use homography to map image points to BEV coordinates. Learn the Roboflow keypoint and bounding-box annotation procedures and prepare a small dataset for Day 2.}{A working Colab notebook, image points and matching BEV points, a homography matrix, a single-point projection, and confirmed keypoint and box annotation formats.}

\vspace{1mm}
\footnotesize
\begin{xltabular}{\textwidth}{|C{0.85cm}|C{1.2cm}|C{1.55cm}|Y|Y|C{0.75cm}|}\hline
\tasktableheader
\endfirsthead
\hline
\tasktableheader
\endhead
\checkbox & 1.1 & \typebadge{Setup} & Open the course repository and check access to Colab, Google Drive, assets and src. & \checks{\item The notebook displays COURSE\_ROOT and a sample image.} & 3 \\ \hline
\checkbox & 1.2 & \typebadge{Coords} & Use the BEV/camera keypoint tool to define corresponding points. & \checks{\item At least four keypoint pairs are defined.} & 4 \\ \hline
\checkbox & 1.3 & \typebadge{BEV} & Check each BEV keypoint name and its corresponding image location. & \checks{\item Image and BEV point counts match, and each pair's name can be explained.} & 3 \\ \hline
\checkbox & 1.4 & \typebadge{Mapping} & Use cv2.findHomography to estimate the matrix, then project a selected point onto BEV. & \checks{\item Output the matrix and an image showing the projected point.} & 5 \\ \hline
\checkbox & 1.5 & \typebadge{Labels} & Review Roboflow keypoint and box annotation assignments, then label a small set of images. & \checks{\item Confirm the required formats for both assignments.} & 5 \\ \hline
\rowcolor{LightBlue}\multicolumn{5}{|r|}{\textbf{Day 1 subtotal}} & \textbf{20} \\ \hline
\end{xltabular}
\normalsize
\vspace{3mm}
\begin{projectfiles}
\projectfile{day1/d1\_01\_keypoint\_pairing\_homography.ipynb}{Define image-to-BEV point pairs, estimate homography and project a selected point onto BEV.}
\projectfile{day1/d1\_02\_keypoint\_annotation\_roboflow\_lab.ipynb}{Learn keypoint ordering, data formats, error checks and an optional fine-tuning workflow.}
\projectfile{day1/d1\_03\_bbox\_homework.ipynb}{Prepare box annotations for player, ball, rim and number classes, and a Roboflow dataset version for Day 2.}
\end{projectfiles}

\clearpage
\section{Day 2: Detection Previews, YOLO26 and BBOX-to-BEV}
\daybox{Understand detection outputs as collections of bounding boxes before attempting full model training. Start with a manually specified player box and its bottom-center footpoint. Examine sample detections or the instructor's trained YOLO26 model, then project multiple player footpoints onto BEV.}{A single-box BEV projection, detection preview images, an understanding of the training preview, and a multi-player BBOX-to-BEV projection.}

\vspace{1mm}\footnotesize
\begin{xltabular}{\textwidth}{|C{0.85cm}|C{1.2cm}|C{1.55cm}|Y|Y|C{0.75cm}|}\hline
\tasktableheader
\endfirsthead
\hline
\tasktableheader
\endhead
\checkbox & 2.1 & \typebadge{Boxes} & Use a manually specified player box. Take its bottom-center as an estimate of the player's floor position. & \checks{\item Draw the box and its bottom-center point.} & 4 \\ \hline
\checkbox & 2.2 & \typebadge{Project} & Reuse the Day 1 homography to project one player's footpoint onto BEV. & \checks{\item Produce a side-by-side image: box on the left, BEV point on the right.} & 5 \\ \hline
\checkbox & 2.3 & \typebadge{Detect} & Examine player, ball, rim and number detections using sample outputs or a trained YOLO26/Ultralytics-style model. & \checks{\item Output a detection preview image.} & 6 \\ \hline
\checkbox & 2.4 & \typebadge{Data} & Inspect the Roboflow box dataset and training preview. Explain datasets, classes, confidence and model limitations. & \checks{\item Explain one possible cause of detection failure.} & 4 \\ \hline
\checkbox & 2.5 & \typebadge{Combine} & Convert multiple player boxes into footpoints and project them onto BEV. & \checks{\item Output a multi-player BBOX-to-BEV image.} & 6 \\ \hline
\rowcolor{LightBlue}\multicolumn{5}{|r|}{\textbf{Day 2 subtotal}} & \textbf{25} \\ \hline
\end{xltabular}
\normalsize
\vspace{3mm}
\begin{projectfiles}
\projectfile{day2/d2\_01\_yolo26\_detection.ipynb}{Build on Day 1's data preparation and interpret classes, confidence scores, boxes and short-video detections.}
\projectfile{day2/d2\_02\_yolo\_players\_to\_bev.ipynb}{Project bottom-center footpoints from multiple player boxes onto BEV.}
\projectfile{day2/d2\_03\_bbox\_to\_bev\_integration.ipynb}{Combine detections, court keypoints, homography and BEV output frame by frame.}
\end{projectfiles}

\clearpage
\section{Day 3: ByteTrack, Multi-Object Tracking, BEV Paths and Team Clustering}
\daybox{Complete the position-data sequence from image coordinates and homography to detection and tracking. Use IoU to understand box association across frames, inspect ByteTrack IDs and paths, and project each track onto BEV. Finally, use jersey torso crops, HSV histograms and K-means to form two teams.}{An IoU heatmap with association lines, video with track IDs and paths, BEV path data, and Team A/B clustering visualizations.}

\vspace{1mm}\footnotesize
\begin{xltabular}{\textwidth}{|C{0.85cm}|C{1.2cm}|C{1.55cm}|Y|Y|C{0.75cm}|}\hline
\tasktableheader
\endfirsthead
\hline
\tasktableheader
\endhead
\checkbox & 3.1 & \typebadge{Concept} & Compute IoU between boxes in two consecutive frames to understand box association. & \checks{\item Output an IoU matrix and explain the matches.} & 4 \\ \hline
\checkbox & 3.2 & \typebadge{Track} & Apply ByteTrack principles or a simplified tracker to assign IDs to detections across frames. & \checks{\item Output a frame annotated with track IDs.} & 5 \\ \hline
\checkbox & 3.3 & \typebadge{Paths} & Collect box-center or footpoint coordinates for each track ID to form player paths. & \checks{\item Export tracks as JSON or a DataFrame.} & 4 \\ \hline
\checkbox & 3.4 & \typebadge{BEV} & Project tracked footpoints onto BEV and draw each ID's path. & \checks{\item Output a BEV trajectory image.} & 5 \\ \hline
\checkbox & 3.5 & \typebadge{Cluster} & Crop players' upper bodies and use HSV features with K-means to form Team A/B. Discuss how track IDs differ from real player identities. & \checks{\item Compare boxes, torso crops, 2D feature plots and team labels; identify one clustering limitation.} & 2 \\ \hline
\rowcolor{LightBlue}\multicolumn{5}{|r|}{\textbf{Day 3 subtotal}} & \textbf{20} \\ \hline
\end{xltabular}
\normalsize
\vspace{3mm}
\begin{projectfiles}
\projectfile{day3/d3\_01\_tracking\_concept\_iou\_association.ipynb}{Use an IoU matrix to understand how boxes are matched across frames.}
\projectfile{day3/d3\_02\_yolo\_bytetrack\_tracking.ipynb}{Use YOLO and ByteTrack to examine track IDs, short paths and track lifecycle summaries.}
\projectfile{day3/d3\_03\_team\_clustering.ipynb}{Visualize two-team clustering with torso crops, HSV histograms and two-cluster K-means.}
\projectfile{day3/d3\_04\_tracking\_to\_bev\_mini\_project.ipynb}{Combine track IDs, BEV paths and track-level team clustering. Output team-colored boxes on the original image, a feature scatterplot and team-colored BEV paths.}
\end{projectfiles}

\clearpage
\section{Day 4: Close-Up Shooting Video, Human Pose and Ball Trajectories}
\daybox{Begin the second part of the course with close-up shooting footage. Record a shot on a phone, place the video in Google Drive's assets/raw folder or select it in the Day 4-01 Colab upload cell, then convert it to MP4. Use the provided shot detector, YOLO pose and ball detector. Combine shot-candidate scores, wrist positions and ball trajectories to estimate the release frame, joint angles and 2D launch angle.}{A converted shooting video, shot-event CSV, pose-angle CSV, event-specific ball tracks, angle plots, release-check images and an overlay video.}

\vspace{1mm}\footnotesize
\begin{xltabular}{\textwidth}{|C{0.85cm}|C{1.2cm}|C{1.55cm}|Y|Y|C{0.75cm}|}\hline
\tasktableheader
\endfirsthead
\hline
\tasktableheader
\endhead
\checkbox & 4.1 & \typebadge{Upload} & Place the video in Google Drive's assets/raw folder, or set USE\_BROWSER\_UPLOAD=True in Day 4-01 and select a file. & \checks{\item The ``using video'' output points to the expected converted video.} & 3 \\ \hline
\checkbox & 4.2 & \typebadge{Convert} & Use ffmpeg to produce MP4 at 30 FPS, with a maximum side length of 1280 pixels and no audio. & \checks{\item An MP4 file appears in assets/converted.} & 3 \\ \hline
\checkbox & 4.3 & \typebadge{Events} & Use the shot detector and YOLO pose to identify the shooter and candidate shot intervals. & \checks{\item Output frame-level shot scores and a pose DataFrame.} & 3 \\ \hline
\checkbox & 4.4 & \typebadge{Angles} & Compute elbow, knee and shoulder angles, then read the shooting-side angles at the release frame. & \checks{\item Output angle curves and joint angles at release.} & 3 \\ \hline
\checkbox & 4.5 & \typebadge{Ball path} & Initialize an event-specific ball track near the shooter's wrist. Estimate the release frame and 2D launch angle. & \checks{\item Output an event CSV, launch angle and images for manual review.} & 3 \\ \hline
\rowcolor{LightBlue}\multicolumn{5}{|r|}{\textbf{Day 4 subtotal}} & \textbf{15} \\ \hline
\end{xltabular}
\normalsize
\vspace{3mm}
\begin{projectfiles}
\projectfile{day4/d4\_01\_ball\_detector\_bytetrack\_preview.ipynb}{Prepare your video and inspect ball boxes and track IDs with the provided ball detector and ByteTrack.}
\projectfile{day4/d4\_02\_yolo\_pose\_and\_shot\_lab.ipynb}{Use the shot detector and YOLO pose to obtain candidate shot scores and the shooter's joint angles.}
\projectfile{day4/d4\_03\_ball\_tracking\_and\_release\_point\_lab.ipynb}{Combine event, wrist and ball-track signals to estimate the release frame, joint angles and 2D launch angle, with visual checks.}
\end{projectfiles}

\clearpage
\section{Day 5: Project Proposals and Completed Project Demonstrations}
\daybox{No new prescribed implementation is introduced. Apply Days 1--4 to propose a feasible, evaluable basketball vision project or demonstrate a completed project. Explain the problem, data and methods, then present evidence, limitations and next steps.}{6--10 slides, a one-page summary, a system diagram and an evaluation plan. Proposals also include milestones and expected results; completed demonstrations include reproducible outputs.}

\vspace{1mm}\footnotesize
\begin{xltabular}{\textwidth}{|C{0.85cm}|C{1.2cm}|C{1.55cm}|Y|Y|C{0.75cm}|}\hline
\tasktableheader
\endfirsthead
\hline
\tasktableheader
\endhead
\checkbox & 5.1 & \typebadge{Problem} & Define users, inputs, outputs and one to three answerable research or product questions. & \checks{\item State a clear problem rather than a list of models.} & 4 \\ \hline
\checkbox & 5.2 & \typebadge{Data} & Explain data sources, annotations, splits, permissions and assumed prerequisite modules. & \checks{\item Distinguish completed work, reused tools and planned work.} & 3 \\ \hline
\checkbox & 5.3 & \typebadge{Method} & Diagram the flow from video through detections, tracks, team/identity and BEV/pose/events to the final output. & \checks{\item Explain each module's inputs, outputs and main assumptions.} & 4 \\ \hline
\checkbox & 5.4 & \typebadge{Evidence} & Proposals specify baselines, metrics and expected charts. Completed demonstrations include representative successes and failures. & \checks{\item The evaluation addresses the original question.} & 4 \\ \hline
\checkbox & 5.5 & \typebadge{Present} & Present the proposal or demo within the allotted time, using charts, video or a storyboard to support conclusions. & \checks{\item Complete slides, summary and required attachments.} & 3 \\ \hline
\checkbox & 5.6 & \typebadge{Review} & Explain limitations, risks, error propagation and prioritized next steps. & \checks{\item Propose at least one feasible follow-up experiment.} & 2 \\ \hline
\rowcolor{LightBlue}\multicolumn{5}{|r|}{\textbf{Day 5 subtotal}} & \textbf{20} \\ \hline
\end{xltabular}
\normalsize
\vspace{3mm}
\begin{projectfiles}
\projectfile{day5/project\_proposal\_spec.tex}{Submission formats, presentation timing, assessment criteria and suggested topics for proposals and completed demonstrations.}
\projectfile{daily\_lecture\_outline.tex}{Outline for the instructor's opening lecture hour each day; Day 5 covers research design, evaluation and technical communication.}
\end{projectfiles}
\par\small\textbf{Source-file note:} The instructor file \texttt{daily\_lecture\_outline.tex} is listed in the original schedule but is not included in this repository. For the theory, use the full course handout listed in the README.\normalsize


\clearpage
\section{Bonus Tasks}
\footnotesize
\begin{xltabular}{\textwidth}{|C{0.85cm}|C{1.2cm}|C{1.55cm}|Y|Y|C{0.75cm}|}\hline
\tasktableheader
\endfirsthead
\hline
\tasktableheader
\endhead
\checkbox & E.1 & \typebadge{Bonus} & Improve a demo: adjust homography points, tune the ball threshold, analyze your own video, or improve output figures. & \checks{\item Show an additional modification and its result.} & +5 \\ \hline
\checkbox & E.2 & \typebadge{Feedback} & Submit specific feedback on tools, task difficulty, materials or possible extensions. & \checks{\item Provide concrete feedback or suggestions.} & +5 \\ \hline
\rowcolor{WarnYellow}\multicolumn{5}{|r|}{\textbf{Bonus total}} & \textbf{+10} \\ \hline
\end{xltabular}
\normalsize

\end{document}
